\chapter{Frames, Rotations and Quaternions}
\label{ch:rotations}

\begin{objectives}
\begin{itemize}
  \item Read and write frame-annotated transformations such as $R_{wb}$, $p_{wb}$ and $T_{bc}$ without
        ambiguity, and catch frame bugs with the adjacency rule.
  \item Derive Rodrigues' formula and understand the axis-angle (rotation vector) representation.
  \item See exactly why Euler angles are banned from the estimator.
  \item Work fluently with Hamilton unit quaternions, and avoid the Eigen and EuRoC ordering traps.
  \item Derive the kinematic equations $\dot R = R\sx{\omega}$ and $\dot q = \tfrac12 q\qprod[0,\omega]$.
\end{itemize}
\end{objectives}

% ------------------------------------------------------------------------------
\section{Frames and notation}

A \emph{frame} is an origin together with three orthonormal, right-handed axes. Three frames appear
in this project (\cref{fig:frames}).

\begin{convention}[frames, AGENTS.md C1]
\begin{itemize}
  \item $W$, the \textbf{world} frame: fixed, gravity-aligned with $z$ pointing up, the frame of the
        ground truth (EuRoC calls it \code{R}). Gravity is $g_w = [0, 0, -9.81]\T~\mathrm{m/s^2}$.
  \item $B$, the \textbf{body} frame: the IMU frame \code{imu0} (EuRoC's \code{S} in \code{imu0}).
  \item $C$, the left \textbf{camera} frame \code{cam0}: origin at the optical center, $z$ along the
        optical axis, $x$ to the right and $y$ down in the image.
\end{itemize}
$R_{wb}$ (or $q_{wb}$) maps body-frame coordinates to world-frame coordinates, $v_w = R_{wb}v_b$.
$p_{wb}$ is the position of the origin of $B$ expressed in $W$. $v_{wb}$ is the velocity of $B$
relative to $W$, expressed in $W$. $T_{bc} = (R_{bc}, p_{bc})$ maps camera points into the body frame,
$p_b = R_{bc}p_c + p_{bc}$; it is \code{T\_BS} from \code{cam0/sensor.yaml}. A world-to-body rotation is
never used as a state.
\end{convention}

\begin{figure}[t]
\centering
\begin{tikzpicture}[scale=1.0]
  % world frame
  \coordinate (W) at (0,0);
  \draw[frameaxis, plotorange] (W) -- ++(1.4,0) node[right] {$x_w$};
  \draw[frameaxis, plotgreen] (W) -- ++(0.75,0.6) node[above right] {$y_w$};
  \draw[frameaxis, plotblue] (W) -- ++(0,1.4) node[above] {$z_w$};
  \node[below left] at (W) {$W$};
  \draw[->, thick, black!60] (-0.6,1.0) -- ++(0,-0.8) node[midway, left] {$g$};
  % body frame
  \coordinate (B) at (5.2,2.2);
  \begin{scope}[shift={(B)}, rotate=25]
    \draw[fill=primarylight, draw=primary, rounded corners=2pt] (-0.55,-0.25) rectangle (0.55,0.25);
    \draw[frameaxis, plotorange] (0,0) -- (1.3,0) node[right] {$x_b$};
    \draw[frameaxis, plotgreen] (0,0) -- (0.65,0.55) node[above right] {$y_b$};
    \draw[frameaxis, plotblue] (0,0) -- (0,1.3) node[above] {$z_b$};
    \node[below left] at (-0.2,-0.25) {$B$ (IMU)};
    % camera
    \coordinate (C) at (2.0,-1.3);
    \draw[fill=white, draw=primary] ($(C)+(-0.15,0.15)$) rectangle ($(C)+(0.15,-0.15)$);
    \draw[frameaxis, plotblue] (C) -- ++(0,-0.9) node[below] {$z_c$};
    \draw[frameaxis, plotorange] (C) -- ++(0.8,0) node[right] {$x_c$};
    \draw[frameaxis, plotgreen] (C) -- ++(-0.45,-0.35) node[left] {$y_c$};
    \node[above right] at ($(C)+(0.1,0.1)$) {$C$};
    \draw[->, thick, conv] (0.3,-0.3) to[bend right=20] node[midway, below left=-1pt] {$T_{bc}$} ($(C)+(-0.2,0.05)$);
  \end{scope}
  \draw[->, thick, primary, dashed] (W) -- (B) node[midway, above left] {$p_{wb}$, $R_{wb}$};
\end{tikzpicture}
\caption{The world frame $W$ (gravity-aligned, $z$ up), the body/IMU frame $B$ moving with the
vehicle, and the camera frame $C$ rigidly attached to the body through the constant extrinsic
calibration $T_{bc}$.}
\label{fig:frames}
\end{figure}

\subsection{Reading subscripts}

The naming rule \code{<quantity>\_<to><from>} makes products self-checking. $R_{wb}$ takes a vector
\emph{from} $B$ \emph{to} $W$. Composition chains frames like dominoes:
\begin{equation}
  R_{wc} = R_{wb}\,R_{bc}, \qquad v_w = R_{wb}\,R_{bc}\,v_c .
\end{equation}

\begin{keyidea}[the adjacency rule]
In a valid product, adjacent subscripts match: $R_{w\underline{b}}\,R_{\underline{b}c}$. A product such as
$R_{wb}R_{cb}$ has mismatched neighbors and is a bug. The fix is a transpose,
$R_{wb}R_{cb}\T = R_{wb}R_{bc}$, since $R_{bc} = R_{cb}\inv = R_{cb}\T$. Apply the rule to every line of
code that multiplies rotations or transforms points.
\end{keyidea}

\subsection{Active and passive views}

A rotation matrix can be read in two ways. In the \emph{passive} view, $R_{wb}$ re-expresses the
same physical vector in another frame: $v_w$ and $v_b$ are coordinates of one arrow. In the
\emph{active} view, $R_{wb}$ rotates the world axes onto the body axes: its columns are the body
axes $x_b, y_b, z_b$ expressed in world coordinates. Both readings are the same matrix:
\begin{equation}
  R_{wb} = \begin{bmatrix} x_b^{(w)} & y_b^{(w)} & z_b^{(w)} \end{bmatrix}.
\end{equation}
This gives a quick sanity check when debugging. The third column of $R_{wb}$ is the body $z$-axis
in world coordinates. For a level vehicle it is $[0,0,1]\T$.

% ------------------------------------------------------------------------------
\section{Rotation matrices and \texorpdfstring{$\SO$}{SO(3)}}

\begin{definition}[Special orthogonal group]
$\SO = \{ R \in \R^{3\times3} : R\T R = I,\ \det R = +1 \}$.
\end{definition}

Orthogonality ($R\T R = I$) means $R$ preserves lengths and angles. $\det R = +1$ excludes
reflections, which have $\det = -1$ and would turn a right-handed frame into a left-handed one. The
set is closed under multiplication and inversion, $R\inv = R\T$, so it is a \emph{group}. A rotation
matrix has nine entries but only three degrees of freedom, since orthonormality imposes six
constraints. This redundancy is why $\SO$ is a curved three-dimensional surface inside $\R^9$, not a
vector space. The sum of two rotation matrices, for example, is not a rotation matrix.

The elementary rotations about the coordinate axes are
\begin{equation}
  R_x(\alpha) = \begin{bmatrix} 1 & 0 & 0 \\ 0 & c_\alpha & -s_\alpha \\ 0 & s_\alpha & c_\alpha \end{bmatrix},\quad
  R_y(\alpha) = \begin{bmatrix} c_\alpha & 0 & s_\alpha \\ 0 & 1 & 0 \\ -s_\alpha & 0 & c_\alpha \end{bmatrix},\quad
  R_z(\alpha) = \begin{bmatrix} c_\alpha & -s_\alpha & 0 \\ s_\alpha & c_\alpha & 0 \\ 0 & 0 & 1 \end{bmatrix},
\end{equation}
with $c_\alpha = \cos\alpha$ and $s_\alpha = \sin\alpha$. Matrix multiplication of rotations does not
commute: $R_x(90^\circ)R_y(90^\circ) \ne R_y(90^\circ)R_x(90^\circ)$. Try it with a book.

% ------------------------------------------------------------------------------
\section{Axis-angle and Rodrigues' formula}

\begin{theorem}[Euler's rotation theorem]
Every $R\in\SO$ is a rotation by some angle $\theta \in [0, \pi]$ about some unit axis $u$, the
eigenvector of $R$ with eigenvalue $1$.
\end{theorem}

The pair $(u, \theta)$ is the \emph{axis-angle} representation, and $\phi = \theta u \in \R^3$ is the
\emph{rotation vector}. Its direction is the axis and its length is the angle. The rotation vector is
the central object of the next chapter.

\begin{theorem}[Rodrigues' rotation formula]
\label{thm:rodrigues}
The rotation by angle $\theta$ about the unit axis $u$ is
\begin{equation}
  R = I + \sin\theta\,\sx{u} + (1 - \cos\theta)\,\sx{u}^2
    = \cos\theta\, I + (1-\cos\theta)\,u u\T + \sin\theta\,\sx{u}.
  \label{eq:rodrigues}
\end{equation}
\end{theorem}
\begin{proof}[Geometric proof]
Split any vector into components parallel and perpendicular to the axis,
$v = v_\parallel + v_\perp$ with $v_\parallel = u u\T v$. Rotation leaves $v_\parallel$ unchanged and
rotates $v_\perp$ by $\theta$ within the plane spanned by $v_\perp$ and $u \times v_\perp = u \times v$:
\begin{equation*}
  Rv = v_\parallel + \cos\theta\, v_\perp + \sin\theta\, u \times v .
\end{equation*}
Substituting $v_\perp = v - uu\T v = -\sx{u}^2 v$ (from \eqref{eq:skew-sq} with $\|u\|=1$) and
$u \times v = \sx{u}v$ gives $Rv = (I + \sin\theta\sx{u} + (1 - \cos\theta)\sx{u}^2)v$.
\end{proof}
\begin{proof}[Algebraic proof]
The rotation about $u$ at unit angular rate satisfies $\dot R = \sx{u} R$ (\cref{sec:rot-kinematics}),
so $R(\theta) = \exp(\theta\sx{u})$. With $\sx{u}^3 = -\sx{u}$ from \eqref{eq:skew-cube}, the
exponential series regroups as
\begin{equation*}
  \exp(\theta\sx{u}) = I + \Big(\theta - \tfrac{\theta^3}{3!} + \tfrac{\theta^5}{5!} - \cdots\Big)\sx{u}
     + \Big(\tfrac{\theta^2}{2!} - \tfrac{\theta^4}{4!} + \cdots\Big)\sx{u}^2,
\end{equation*}
and the two series are $\sin\theta$ and $1 - \cos\theta$.
\end{proof}

Two consequences of \eqref{eq:rodrigues} let us go back from $R$ to $(u, \theta)$. Taking the trace
(using $\tr\sx{u} = 0$ and $\tr\sx{u}^2 = -2$) and the skew-symmetric part,
\begin{equation}
  \tr R = 1 + 2\cos\theta, \qquad R - R\T = 2\sin\theta\,\sx{u}.
  \label{eq:trace-angle}
\end{equation}
\Cref{ch:lie} shows that \emph{how} these are combined numerically matters a great deal.

% ------------------------------------------------------------------------------
\section{Euler angles and gimbal lock}

Euler angles describe a rotation as three successive rotations about coordinate axes. The aerospace
roll-pitch-yaw ($ZYX$) convention writes
\begin{equation}
  R_{wb} = R_z(\psi)\,R_y(\vartheta)\,R_x(\varphi),
\end{equation}
with yaw $\psi$, pitch $\vartheta$ and roll $\varphi$. They are intuitive, which is why they are
used for display. Their fatal flaw appears in the kinematics. The relation between the body angular
velocity measured by a gyroscope and the Euler angle rates is
\begin{equation}
  \begin{bmatrix} \dot\varphi \\ \dot\vartheta \\ \dot\psi \end{bmatrix}
  = \begin{bmatrix}
      1 & \sin\varphi\tan\vartheta & \cos\varphi\tan\vartheta \\
      0 & \cos\varphi & -\sin\varphi \\
      0 & \sin\varphi/\cos\vartheta & \cos\varphi/\cos\vartheta
    \end{bmatrix} \omega_b .
  \label{eq:euler-rates}
\end{equation}
At $\vartheta = \pm90^\circ$, $\cos\vartheta = 0$ and the matrix blows up. Geometrically, when the
nose points straight up, yaw and roll rotate about the same physical axis, so the three angles no
longer describe three independent motions. This is \emph{gimbal lock}, named after mechanical
gimbal platforms that physically jam in this configuration. Near the singularity, small angular
velocities produce enormous Euler-angle rates, and a filter that linearizes in Euler angles becomes
numerically meaningless.

\begin{convention}[AGENTS.md C2]
Euler angles never appear in the estimator. They may appear only in telemetry and plots.
\end{convention}

No three-parameter representation of $\SO$ is free of singularities everywhere, a topological fact
(the space $\SO$ cannot be covered by a single chart). The solution adopted by modern estimators is to
store attitude in a singularity-free \emph{redundant} representation (a unit quaternion or a rotation
matrix) and to use three parameters only for \emph{small} deviations around the current estimate, where
no singularity is near. That is the content of the next chapter.

% ------------------------------------------------------------------------------
\section{Quaternions}
\label{sec:quaternions}

\subsection{Quaternion algebra}

A quaternion is a four-component number $q = w + x\,i + y\,j + z\,k$, which we write as a scalar part
and a vector part, $q = [w,\ v]$ with $v = (x,y,z)\T$. Hamilton's multiplication rules are
\begin{equation}
  i^2 = j^2 = k^2 = ijk = -1, \qquad\text{hence}\qquad ij = k,\ jk = i,\ ki = j,\ ji = -k, \dots
\end{equation}
Expanding the product of two quaternions with these rules gives
\begin{equation}
  p \qprod q = \begin{bmatrix} p_w q_w - p_v\T q_v \\ p_w q_v + q_w p_v + p_v \times q_v \end{bmatrix}.
  \label{eq:qprod}
\end{equation}
The cross product makes the product non-commutative. Because it is linear in each factor, it can be
written with matrices,
\begin{equation}
  p\qprod q = [p]_L\, q = [q]_R\, p, \qquad
  [p]_L = \begin{bmatrix} p_w & -p_v\T \\ p_v & p_w I + \sx{p_v} \end{bmatrix}, \quad
  [q]_R = \begin{bmatrix} q_w & -q_v\T \\ q_v & q_w I - \sx{q_v} \end{bmatrix},
\end{equation}
where $q$ on the right of $[p]_L\,q$ is the 4-vector $[q_w; q_v]$. The \emph{conjugate} is
$q^* = [w,\ -v]$, the norm is $\|q\| = \sqrt{q\qprod q^*} = \sqrt{w^2 + \|v\|^2}$, and
$q\inv = q^*/\|q\|^2$. Products reverse under conjugation: $(p\qprod q)^* = q^*\qprod p^*$.

\subsection{Unit quaternions represent rotations}

\begin{theorem}[Rotation by a unit quaternion]
\label{thm:quat-rotation}
For a unit quaternion $q = [w, v]$, the map $a \mapsto q\qprod[0, a]\qprod q^*$ sends pure
quaternions to pure quaternions and acts on $a$ as the rotation matrix
\begin{equation}
  R(q) = (w^2 - v\T v)\,I + 2\,v v\T + 2w\,\sx{v}
  = \begin{bmatrix}
    1 - 2(y^2 + z^2) & 2(xy - wz) & 2(xz + wy) \\
    2(xy + wz) & 1 - 2(x^2 + z^2) & 2(yz - wx) \\
    2(xz - wy) & 2(yz + wx) & 1 - 2(x^2 + y^2)
  \end{bmatrix}.
  \label{eq:quat-to-rot}
\end{equation}
The unit quaternion of a rotation by $\theta$ about the unit axis $u$ is
\begin{equation}
  q = \big[\cos(\theta/2),\ \sin(\theta/2)\,u\big].
  \label{eq:quat-axis-angle}
\end{equation}
\end{theorem}
\begin{proof}
Expand $q\qprod[0, a]\qprod q^*$ twice with \eqref{eq:qprod}. The scalar part is $0$ and the vector part
is $(w^2 - v\T v)a + 2(v\T a)v + 2w\, v\times a$, which is \eqref{eq:quat-to-rot}. For
\eqref{eq:quat-axis-angle}, substitute $w = \cos\frac\theta2$, $v = \sin\frac\theta2 u$ and use the
double-angle identities: $w^2 - v\T v = \cos\theta$, $2\sin^2\frac\theta2 = 1 - \cos\theta$ and
$2\sin\frac\theta2\cos\frac\theta2 = \sin\theta$. The result is the second form of Rodrigues'
formula \eqref{eq:rodrigues}.
\end{proof}

The \emph{half} angle in \eqref{eq:quat-axis-angle} is the most common source of factor-of-two bugs. It
arises because $q$ appears twice in the sandwich product $q \qprod [0,a] \qprod q^*$.

\begin{keyidea}[double cover]
$q$ and $-q$ give the same $R(q)$, since \eqref{eq:quat-to-rot} is quadratic in $q$. Every rotation
has exactly two unit quaternions, at opposite points of the unit sphere $S^3 \subset \R^4$. They
correspond to rotating by $\theta$ about $u$, or by $2\pi - \theta$ about $-u$. Code that compares,
averages, interpolates or takes logarithms of quaternions must account for this. The project
canonicalizes to $w \ge 0$ before taking $\Log_q$ (AGENTS.md C2).
\end{keyidea}

\subsection{Composition and the Hamilton convention}

With the Hamilton product, quaternion multiplication composes rotations in the same order as matrix
multiplication:
\begin{equation}
  R(p \qprod q) = R(p)\,R(q), \qquad q_{wc} = q_{wb}\qprod q_{bc} \ \leftrightarrow\ R_{wc} = R_{wb}R_{bc}.
\end{equation}
The adjacency rule therefore applies to quaternions unchanged.

Unfortunately, two incompatible quaternion conventions are in wide use. Much of the classic VIO
literature, including the original MSCKF paper~\cite{mourikis2007msckf} and the indirect Kalman filter
report of Trawny and Roumeliotis~\cite{trawny2005indirect}, uses the \emph{JPL} convention. Formulas
cannot be copied between the two without translation~\cite{sommer2018flipped}.

\begin{table}[h]
\centering\small
\begin{tabularx}{\textwidth}{@{}l X X@{}}
\toprule
 & \textbf{Hamilton (this project)} & \textbf{JPL} \\
\midrule
Multiplication rule & $ij = k$ ($ijk = -1$) & $ij = -k$ ($ijk = +1$) \\
Component order & usually $[w, x, y, z]$ & usually $[x, y, z, w]$ \\
Typical meaning & body-to-world (active), $q_{wb}$ & world-to-body (passive), ${}^{B}_{G}\bar q$ \\
Composition & $q_{wc} = q_{wb}\qprod q_{bc}$ & ${}^{C}_{G}\bar q = {}^{C}_{B}\bar q\otimes{}^{B}_{G}\bar q$ \\
Used by & Eigen, ROS, Solà~\cite{sola2017quaternion}, Ceres, Sophus & MSCKF papers, Trawny~\cite{trawny2005indirect}, NASA \\
\bottomrule
\end{tabularx}
\caption{Hamilton versus JPL quaternions. The JPL product of two quaternions equals the Hamilton
product of the same numbers in reverse order.}
\label{tab:hamilton-jpl}
\end{table}

\begin{convention}[quaternions, AGENTS.md C2]
Hamilton quaternions, \code{Eigen::Quaterniond}, always of unit norm, normalized after every
multiplicative update. Never substitute JPL formulas, even when a paper you are reading uses them.
Translate them first.
\end{convention}

\subsection{Quaternions in code and data}

\begin{pitfall}[coefficient order]
Three orderings meet at the project's input boundary:
\begin{itemize}
  \item \code{Eigen::Quaterniond(w, x, y, z)}: the constructor takes $w$ \emph{first};
  \item \code{q.coeffs()}: Eigen's storage is $[x, y, z, w]$, with $w$ \emph{last};
  \item EuRoC CSV files store \code{qw, qx, qy, qz}, with $w$ first.
\end{itemize}
Mapping a raw EuRoC row with \code{Eigen::Map<Quaterniond>} interprets it as $[x,y,z,w]$ and silently
produces a different, valid-looking rotation. Always construct explicitly,
\code{Eigen::Quaterniond(row[0], row[1], row[2], row[3])}, and cover the parser with a fixture test.
\end{pitfall}

Other Eigen facts worth knowing: \code{q * v} for a \code{Vector3d} applies $R(q)$ to $v$;
\code{q1 * q2} is the Hamilton product; \code{q.conjugate()} is the inverse of a unit quaternion;
\code{q.normalize()} works in place while \code{q.normalized()} returns a copy; and
\code{q.toRotationMatrix()} evaluates \eqref{eq:quat-to-rot}.

\paragraph{Normalization.} Each multiplicative update of a quaternion rounds its four components,
so $\|q\|$ performs a slow random walk away from $1$. A non-unit quaternion inserted into
\eqref{eq:quat-to-rot} produces a matrix that scales vectors as well as rotating them. Dividing by
the norm after every update costs a few flops and removes the problem. Milestone 4 checks that
$|\|q\| - 1| < 10^{-12}$ after $10^6$ updates.

\paragraph{From a rotation matrix to a quaternion.} Inverting \eqref{eq:quat-to-rot} naively via
$w = \tfrac12\sqrt{1 + \tr R}$ and $v = \frac{(R - R\T)\veeop}{4w}$ divides by a number that vanishes
for $180^\circ$ rotations. Robust implementations (Shepperd's method, used by Eigen) first determine
which of $w, x, y, z$ has the largest magnitude from the diagonal of $R$, compute that one with a
square root, and obtain the other three by division by it.

\paragraph{Interpolation.} Ground truth and IMU data rarely share timestamps with images. The correct
way to interpolate between two attitudes is along the shortest great arc (SLERP), which in the notation
of the next chapter reads
$q(\tau) = q_0 \qprod \Exp_q\!\big(\tau\,\Log_q(q_0\inv\qprod q_1)\big)$ for $\tau\in[0,1]$.
Component-wise linear interpolation followed by normalization is an acceptable approximation only for
small angles.

% ------------------------------------------------------------------------------
\section{Rigid-body transformations}

A rigid transformation combines a rotation and a translation. Point coordinates transform as
\begin{equation}
  p_w = R_{wb}\,p_b + p_{wb}
  \quad\Longleftrightarrow\quad
  \begin{bmatrix} p_w \\ 1 \end{bmatrix} = \underbrace{\begin{bmatrix} R_{wb} & p_{wb} \\ 0 & 1 \end{bmatrix}}_{T_{wb}}
  \begin{bmatrix} p_b \\ 1 \end{bmatrix}.
\end{equation}
The $4\times4$ homogeneous matrices form the group $\SE$. Composition and inversion follow from block
multiplication:
\begin{equation}
  T_{wc} = T_{wb}T_{bc} = \begin{bmatrix} R_{wb}R_{bc} & R_{wb}p_{bc} + p_{wb} \\ 0 & 1\end{bmatrix},
  \qquad
  T_{wb}\inv = T_{bw} = \begin{bmatrix} R_{wb}\T & -R_{wb}\T p_{wb} \\ 0 & 1 \end{bmatrix}.
  \label{eq:se3-compose}
\end{equation}
The camera pose in the world therefore is $R_{wc} = R_{wb}R_{bc}$ and $p_{wc} = p_{wb} + R_{wb}p_{bc}$, and
a world point $p_f$ has camera coordinates
\begin{equation}
  p_c = R_{bc}\T\big(R_{wb}\T(p_f - p_{wb}) - p_{bc}\big),
  \label{eq:world-to-camera}
\end{equation}
which is the first half of the camera measurement model (AGENTS.md C6).

\begin{pitfall}[points versus directions]
Translations apply to points, not to free vectors. A velocity, a gravity vector or a bearing direction
transforms with the rotation only, $v_w = R_{wb}v_b$. In homogeneous coordinates directions carry a $0$
in the fourth component, which drops the translation automatically.
\end{pitfall}

% ------------------------------------------------------------------------------
\section{Rotational kinematics}
\label{sec:rot-kinematics}

\subsection{Angular velocity}

Let $R(t) = R_{wb}(t)$ describe a rotating body. Differentiating $R\T R = I$ gives
$\dot R\T R + R\T\dot R = 0$, so $R\T\dot R$ is skew-symmetric and equals $\sx{\omega_b}$ for a unique
vector $\omega_b$:
\begin{equation}
  \dot R_{wb} = R_{wb}\,\sx{\omega_b} = \sx{\omega_w}\,R_{wb}, \qquad \omega_w = R_{wb}\,\omega_b ,
  \label{eq:rdot}
\end{equation}
where the second form uses \eqref{eq:skew-rot}. $\omega_b$ is the angular velocity of $B$ relative to
$W$, expressed in $B$, and it is exactly what a strapdown gyroscope measures (up to bias and noise).
The two forms of \eqref{eq:rdot} are the kinematic origin of the two perturbation conventions of the
next chapter. A body-frame quantity multiplies on the right, and a world-frame quantity on the left.

\begin{proposition}[Constant body rate]
\label{prop:const-rate}
If $\omega_b$ is constant on $[t_0, t_0 + \Dt]$, then $R(t_0 + \Dt) = R(t_0)\exp(\Dt\sx{\omega_b})$.
\end{proposition}
\begin{proof}
$R(t) = R(t_0)\exp((t - t_0)\sx{\omega_b})$ satisfies \eqref{eq:rdot} by \cref{prop:matexp}, and the
solution is unique.
\end{proof}

This exact solution is the basis of the multiplicative attitude update used in every integration
scheme of \cref{ch:strapdown}. The body rate is held constant over the step and the incremental
rotation is applied on the right. Note that the solution is \emph{not}
$R(t_0) + \Dt\, R(t_0)\sx{\omega_b}$, which would leave $\SO$.

\subsection{Quaternion kinematics}

The quaternion counterpart of \eqref{eq:rdot} is
\begin{equation}
  \dot q_{wb} = \tfrac12\, q_{wb} \qprod [0,\ \omega_b],
  \label{eq:qdot}
\end{equation}
with constant-rate solution $q(t_0 + \Dt) = q(t_0)\qprod\big[\cos\tfrac{\|\omega\|\Dt}{2},\
\sin\tfrac{\|\omega\|\Dt}{2}\,\tfrac{\omega}{\|\omega\|}\big]$. To derive \eqref{eq:qdot}, compose
$q(t + \dd t) = q(t) \qprod \dd q$, where $\dd q$ is the small body-frame rotation by
$\omega_b\,\dd t$. By \eqref{eq:quat-axis-angle}, $\dd q = [1,\ \tfrac12\omega_b\,\dd t] + O(\dd t^2)$,
so $q(t + \dd t) - q(t) = q\qprod[0, \tfrac12\omega_b]\,\dd t$. The factor $\tfrac12$ is the half angle
again.

\subsection{Derivatives of vectors in a rotating frame}

Differentiating $v_w = R_{wb}v_b$ with \eqref{eq:rdot} gives the \emph{transport theorem}
\begin{equation}
  \dot v_w = R_{wb}\big(\dot v_b + \omega_b \times v_b\big).
\end{equation}
It is the source of Coriolis and centripetal terms when dynamics are written in a rotating frame.
Our estimator keeps position and velocity in the (approximately inertial) world frame, so these terms
do not appear in the state equations.

\begin{remark}[Earth rotation]
The world frame is fixed to the Earth, which rotates at $7.29\times10^{-5}$~rad/s. A gyroscope senses
this rate on top of the vehicle's motion. For the ADIS16448 in EuRoC this is comparable to the bias
instability and well below the noise floor over short horizons, so the project neglects it. The
static part of the Earth rate, in the body frame, is partly absorbed by the gyro bias estimate. High-grade
inertial navigation systems model it explicitly.
\end{remark}

\begin{takeaways}
\begin{itemize}
  \item $R_{wb}$ maps body coordinates to world coordinates; its columns are the body axes in the world.
        Adjacent subscripts must match.
  \item Every rotation is an axis-angle rotation, $R = I + \sin\theta\sx{u} + (1-\cos\theta)\sx{u}^2$.
  \item Three-parameter representations have singularities. Euler angles stay out of the estimator.
  \item Hamilton unit quaternions compose like matrices. $q$ and $-q$ are the same rotation, and the
        quaternion of a rotation by $\theta$ uses $\theta/2$.
  \item Eigen's constructor is $w$-first, its storage $w$-last, EuRoC is $w$-first.
  \item $\dot R_{wb} = R_{wb}\sx{\omega_b}$. With constant $\omega_b$ the update
        $R \leftarrow R\exp(\Dt\sx{\omega_b})$ is exact.
\end{itemize}
\end{takeaways}

\begin{exercises}
  \item Show that the columns of any $R \in \SO$ are orthonormal and that $\det R = +1$ implies the
        third column is the cross product of the first two.
  \item Verify $R_x(90^\circ)R_y(90^\circ) \neq R_y(90^\circ)R_x(90^\circ)$ by computing both products,
        and describe each rotation physically.
  \item Derive \eqref{eq:euler-rates} from $R = R_z(\psi)R_y(\vartheta)R_x(\varphi)$ and
        $\sx{\omega_b} = R\T\dot R$.
  \item Prove that $[p]_L[q]_R = [q]_R[p]_L$. What does this say about associativity of the quaternion
        product?
  \item Show that $R(-q) = R(q)$ and compute the quaternion of a $270^\circ$ rotation about $z$ in its
        canonical form ($w \ge 0$). What rotation angle and axis does the canonical form describe?
  \item Write a function that converts a EuRoC ground-truth row into $p_{wb}$, $q_{wb}$, $v_{wb}$, $b_g$
        and $b_a$. Design a fixture row whose values would expose a swapped $w/z$ and swapped biases.
  \item Given $T_{bc}$ from \code{cam0/sensor.yaml} (\cref{app:euroc}), compute the camera's optical axis
        in body coordinates. Which way does the camera point relative to the IMU?
  \item Show that if $\omega_b$ varies in time, $R(t) = R(0)\exp\big(\int_0^t\sx{\omega_b}\,\dd s\big)$ is
        in general \emph{wrong}. Construct a counterexample with piecewise-constant rates (this is the
        ``coning'' error of \cref{ch:strapdown}).
\end{exercises}
